% ═══════════════════════════════════════════════════════════════
% SPANISCH ARBEITSBLATT AB-03c: Los colores + Adjektivangleichung
% ═══════════════════════════════════════════════════════════════
% Sequenz 3: Mi casa | Doppelstunde 3c
% Fachfarbe: #c41e3a (Karminrot)
% Gesamtpunktzahl: 20 Punkte
% ═══════════════════════════════════════════════════════════════

\documentclass[11pt, a4paper]{article}

% SW-MODUS TOGGLE
\newif\ifswmodus
\swmodusfalse

% PAKETE
\usepackage[utf8]{inputenc}
\usepackage[T1]{fontenc}
\usepackage[ngerman]{babel}
\usepackage{helvet}
\renewcommand{\familydefault}{\sfdefault}

\usepackage[margin=2cm]{geometry}
\usepackage{enumitem}
\usepackage{tabularx}
\usepackage{booktabs}
\usepackage{multirow}

\usepackage{xcolor}
\usepackage{framed}
\usepackage{tcolorbox}
\tcbuselibrary{skins}

\usepackage{fancyhdr}
\usepackage{lastpage}
\usepackage{needspace}

% ═══════════════════════════════════════════════════════════════
% FARBEN
% ═══════════════════════════════════════════════════════════════

\definecolor{spanisch-color}{HTML}{c41e3a}
\definecolor{grau-color}{HTML}{888888}
\definecolor{hellgrau-color}{HTML}{f5f5f5}
\definecolor{orange-color}{HTML}{b35c00}

\definecolor{sw-dark}{HTML}{000000}
\definecolor{sw-gray}{HTML}{666666}
\definecolor{sw-white}{HTML}{ffffff}

\ifswmodus
    \colorlet{spanisch}{sw-dark}
    \colorlet{grau}{sw-gray}
    \colorlet{hellgrau}{sw-white}
    \colorlet{orange}{sw-dark}
\else
    \colorlet{spanisch}{spanisch-color}
    \colorlet{grau}{grau-color}
    \colorlet{hellgrau}{hellgrau-color}
    \colorlet{orange}{orange-color}
\fi

\newcommand{\fachfarbe}{spanisch}

% ═══════════════════════════════════════════════════════════════
% VARIABLEN
% ═══════════════════════════════════════════════════════════════

\newcommand{\thema}{Los colores}
\newcommand{\sequenz}{03}
\newcommand{\block}{c}
\newcommand{\fach}{Spanisch}
\newcommand{\klasse}{7}
\newcommand{\maxpunkte}{20}
\newcommand{\datum}{\today}

% ═══════════════════════════════════════════════════════════════
% KOPF- UND FUSSZEILE
% ═══════════════════════════════════════════════════════════════

\pagestyle{fancy}
\fancyhf{}
\renewcommand{\headrulewidth}{0.4pt}
\renewcommand{\footrulewidth}{0.4pt}

\lhead{\fach{} -- Klasse \klasse}
\chead{AB-\sequenz\block}
\rhead{\datum}

\lfoot{}
\cfoot{}
\rfoot{Seite \thepage\ von \pageref{LastPage}}

% ═══════════════════════════════════════════════════════════════
% BEFEHLE
% ═══════════════════════════════════════════════════════════════

\newcommand{\punktebox}[1]{\hfill\fbox{\small \_/#1}}
\newcommand{\luecke}[1]{\rule{#1}{0.4pt}}

\newtcolorbox{merkkasten}[1][]{
    colback=hellgrau,
    colframe=\fachfarbe,
    colbacktitle=white,
    coltitle=\fachfarbe,
    boxrule=1pt,
    arc=0pt,
    left=8pt, right=8pt, top=8pt, bottom=8pt,
    fonttitle=\bfseries,
    attach boxed title to top left={yshift=-2mm, xshift=5mm},
    boxed title style={boxrule=0pt, colback=white},
    title=#1
}

\newtcolorbox{vokabelkasten}{
    colback=white,
    colframe=\fachfarbe,
    boxrule=0.5pt,
    arc=0pt,
    left=5pt, right=5pt, top=5pt, bottom=5pt
}

\newtcolorbox{hinweis}{
    colback=white,
    colframe=orange,
    colbacktitle=white,
    coltitle=orange,
    boxrule=1pt,
    arc=0pt,
    left=5pt, right=5pt, top=5pt, bottom=5pt,
    fonttitle=\bfseries,
    attach boxed title to top left={yshift=-2mm, xshift=5mm},
    boxed title style={boxrule=0pt, colback=white},
    title=Hinweis
}

\newenvironment{aufgabe}[1]{%
    \needspace{5\baselineskip}%
    \nopagebreak[4]%
    \vspace{10pt}
    \noindent\textbf{\textcolor{\fachfarbe}{Aufgabe #1}}
    \nopagebreak[4]%
    \vspace{5pt}
    \nopagebreak[4]%
    \begin{list}{}{\leftmargin=0pt}
    \item[]
}{%
    \end{list}
}

\newcommand{\es}[1]{\textcolor{\fachfarbe}{\textbf{#1}}}

% ═══════════════════════════════════════════════════════════════
% DOKUMENT
% ═══════════════════════════════════════════════════════════════

\begin{document}

% KOPFBEREICH
\begin{center}
    {\Large\bfseries\textcolor{\fachfarbe}{Arbeitsblatt: \thema}}
\end{center}

\vspace{5pt}

\noindent
\begin{tabularx}{\textwidth}{|X|X|X|}
\hline
\textbf{Name:} \luecke{4cm} &
\textbf{Klasse:} \klasse &
\textbf{Datum:} \luecke{2cm} \\
\hline
\end{tabularx}

\vspace{15pt}

% ═══════════════════════════════════════════════════════════════
% MERKKASTEN 1: Die Farben
% ═══════════════════════════════════════════════════════════════

\begin{merkkasten}[Merkkasten 1: Los colores (Die Farben)]
Übertrage die Farben aus dem Lernpfad oder von der Tafel:

\vspace{0.5em}

\begin{tabular}{ll@{\hspace{2cm}}ll}
\es{rojo} & = \luecke{3cm} & \es{blanco} & = \luecke{3cm} \\[0.8em]
\es{azul} & = \luecke{3cm} & \es{marrón} & = \luecke{3cm} \\[0.8em]
\es{verde} & = \luecke{3cm} & \es{gris} & = \luecke{3cm} \\[0.8em]
\es{amarillo} & = \luecke{3cm} & \es{rosa} & = \luecke{3cm} \\[0.8em]
\es{negro} & = \luecke{3cm} & \es{naranja} & = \luecke{3cm} \\
\end{tabular}

\end{merkkasten}

\vspace{10pt}

% ═══════════════════════════════════════════════════════════════
% MERKKASTEN 2: Adjektivangleichung
% ═══════════════════════════════════════════════════════════════

\begin{merkkasten}[Merkkasten 2: Adjektivangleichung]
\textbf{Regeln:}

\vspace{0.5em}

\begin{enumerate}
    \item Adjektive stehen \luecke{2cm} dem Nomen (el sofá \textbf{rojo}).
    \item Adjektive auf \textbf{-o} ändern ihre Endung:
\end{enumerate}

\vspace{0.5em}

\begin{center}
\begin{tabular}{|l|c|c|}
\hline
& \textbf{Singular} & \textbf{Plural} \\
\hline
\textbf{maskulin} & -o (roj\textbf{o}) & \luecke{1.5cm} (roj\textbf{os}) \\
\hline
\textbf{feminin} & \luecke{1.5cm} (roj\textbf{a}) & -as (roj\textbf{as}) \\
\hline
\end{tabular}
\end{center}

\vspace{0.5em}

\begin{enumerate}
    \setcounter{enumi}{2}
    \item Adjektive auf \textbf{-e} oder Konsonant ändern nur den \luecke{2cm}:\\
    verde $\rightarrow$ verdes | azul $\rightarrow$ azules | marrón $\rightarrow$ marrones
\end{enumerate}

\end{merkkasten}

\vspace{15pt}

% ═══════════════════════════════════════════════════════════════
% AUFGABE 1: Farben zuordnen
% ═══════════════════════════════════════════════════════════════

\begin{aufgabe}{1}
Verbinde die spanischen Farben mit den deutschen Übersetzungen. \punktebox{4}
\end{aufgabe}

\begin{center}
\begin{tabular}{rcl}
\es{rojo} & \luecke{3cm} & A) blau \\[0.8em]
\es{azul} & \luecke{3cm} & B) gelb \\[0.8em]
\es{verde} & \luecke{3cm} & C) rot \\[0.8em]
\es{amarillo} & \luecke{3cm} & D) grün \\
\end{tabular}
\end{center}

\vspace{15pt}

% ═══════════════════════════════════════════════════════════════
% AUFGABE 2: Adjektiv anpassen
% ═══════════════════════════════════════════════════════════════

\begin{aufgabe}{2}
Setze das Adjektiv in der richtigen Form ein. \punktebox{4}
\end{aufgabe}

\begin{vokabelkasten}
\textit{Adjektive:} blanco | rojo | azul | negro
\end{vokabelkasten}

\vspace{0.5em}

\noindent
a) La mesa es \luecke{3cm}. (weiß) \\[0.8em]
b) Los sofás son \luecke{3cm}. (rot) \\[0.8em]
c) Las sillas son \luecke{3cm}. (blau) \\[0.8em]
d) El armario es \luecke{3cm}. (schwarz) \\

\vspace{15pt}

% ═══════════════════════════════════════════════════════════════
% AUFGABE 3: Zimmer mit Farben beschreiben
% ═══════════════════════════════════════════════════════════════

\begin{aufgabe}{3}
Beschreibe dein Zimmer. Verwende mindestens 3 Farben. \punktebox{6}
\end{aufgabe}

\begin{hinweis}
Nutze diese Struktur: \es{En mi habitación hay...} oder \es{El/La ... es ...}
\end{hinweis}

\vspace{0.5em}

\noindent
Beispiel: \textit{En mi habitación hay una cama blanca.}

\vspace{0.5em}

\noindent
\rule{\textwidth}{0.4pt}\vspace{8pt}
\rule{\textwidth}{0.4pt}\vspace{8pt}
\rule{\textwidth}{0.4pt}\vspace{8pt}
\rule{\textwidth}{0.4pt}

\vspace{15pt}

% ═══════════════════════════════════════════════════════════════
% AUFGABE 4: Mi casa ideal
% ═══════════════════════════════════════════════════════════════

\begin{aufgabe}{4}
Beschreibe dein Traumhaus (mi casa ideal). Nenne mindestens 3 Zimmer und verwende 4 Farben. \punktebox{6}
\end{aufgabe}

\begin{hinweis}
Beginne mit: \es{Mi casa ideal tiene...} (Mein Traumhaus hat...)\\
Verwende: el salón, la cocina, el dormitorio, el baño, el jardín
\end{hinweis}

\vspace{0.5em}

\noindent
\rule{\textwidth}{0.4pt}\vspace{8pt}
\rule{\textwidth}{0.4pt}\vspace{8pt}
\rule{\textwidth}{0.4pt}\vspace{8pt}
\rule{\textwidth}{0.4pt}\vspace{8pt}
\rule{\textwidth}{0.4pt}\vspace{8pt}
\rule{\textwidth}{0.4pt}

% ═══════════════════════════════════════════════════════════════
% PUNKTEÜBERSICHT
% ═══════════════════════════════════════════════════════════════

\vspace{20pt}
\hrule
\vspace{10pt}

\begin{flushright}
\textbf{Punkte:} \luecke{1cm} / \maxpunkte
\end{flushright}

\end{document}
